\origpage{50}{}
If $k \le w$ then (3) holds, and so we assume that $k = w + 1$. From (2), $|W(e)| \ge w + 1$ for each $e \in E(Q)$, and so $V(W(t_1)) = V(W(t_2))$ if $e \in E(Q)$ has ends $t_1, t_2$, since
\begin{equation*}
W(e) \subseteq V(W(t_1)) \cap V(W(t_2))
\end{equation*}
and $|V(W(t_1))|, |V(W(t_2))| \le w + 1$. Consequently, $V(W(t)) = W(f_1)$ for all $t \in V(Q)$, and in particular $W(f_1) = W(f_2)$. Hence $W'(f) = W(f_1)$ for all $f \in E(P)$, and so $|V(W^*(t))| \le k = w + 1$ for each $t \in V(P) \setminus \{s_1, s_2\}$, by (12.1)(ii) \fix{and the added hypothesis of (12.2)}.\fixnote{fix:12.2-claim3-p50}{Here too, the bound needs that every vertex of $W'(t)$ is in $W'(e)$ for an edge $e$ of $P$ at $t$. Then $V(W'(t)) \subseteq W(f_1)$.} This proves (3).

Let $e_1, e_2 \in E(T^*)$ with $|W^*(e_1)| = |W^*(e_2)| = k'$ say, such that $|W^*(e)| \ge k'$ for all edges $e$ of the path $R$ of $T^*$ with first edge $e_1$ and last edge $e_2$. We must show that there are $k'$ mutually vertex-disjoint paths of $G$ between $W^*(e_1)$ and $W^*(e_2)$. If $e_1, e_2 \in E(P)$ this is clear, and so we may assume that $e_1 \in E(T_1)$. \fix{To see the first case, let $e_1, e_2 \in E(P)$; then $k' = k$ by (12.1)(ii). Let $t \in V(P) \setminus \{s_1, s_2\}$ be incident with the edges $e, e'$ of $P$. By (12.1)(iii) there are $k$ mutually vertex-disjoint paths of $G$ between $W'(e)$ and $W'(e')$: in case (a) the one-vertex paths, and in case (b) the one-vertex paths at the vertices of $X$ and the path formed by the edge of $W'(t)$ with ends $v, v'$. For each edge $e$ of $P$, the unions $C, D$ of the graphs $W'(t)$ over the vertices $t$ on the two sides of $e$ satisfy $C \cup D = G$ and $V(C \cap D) = W'(e)$, since $(P, W')$ is a path-decomposition; so every path of $G$ from $V(C)$ to $V(D)$ meets $W'(e)$. Going along $P$ from $e_1$ to $e_2$ one vertex at a time, and pairing the paths at each $W'(e)$ as in the last part of this proof, we obtain $k$ mutually vertex-disjoint paths of $G$ between $W'(e_1)$ and $W'(e_2)$. This also proves (1), since by (12.1)(i), $W'(e) = W(f_1)$ for the edge $e$ of $P$ incident with $s_1$, and $W'(e) = W(f_2)$ for the edge incident with $s_2$.}\fixnote{fix:12.2-path-proof-added}{The paper gives no argument for this case or for Claim (1). The argument above is the one used in the formalization.} If $e_1, e_2 \in E(T_1)$ the claim follows since $(T, W)$ is linked, and so we may assume that $e_2 \in E(P)$ or $e_2 \in E(T_2)$. If $e_2 \in E(T_2)$ then $k' \le k$ since $P \subseteq R$, and so $|W(e)| \ge k'$ for every edge of $e$ of the path of $T$ between $e_1$ and $e_2$, by (2); and the claim follows since $(T, W)$ is linked. We assume then that $e_2 \in E(P)$, and so $k' = k$. Since $(T, W)$ is linked, there are $k$ mutually disjoint paths of $G$ between $W(e_1)$ and $W(f_1)$, and these are paths of $A = W(T_1)$. Since $(P, W')$ satisfies (12.1), there are $k$ mutually disjoint paths of $G$ between $W(f_1)$ and $W'(e_2)$, and these are paths of $B = W(T_0) \cup W(T_2)$. Since $(A, B)$ is a separation and $V(A \cap B) = W(f_1)$, we can pair these paths to obtain $k$ mutually disjoint paths of $G$ between $W^*(e_1) = W(e_1)$ and $W^*(e_2) = W'(e_2)$, as required. Consequently, $(T^*, W^*)$ is linked. $\square$

We use these lemmas to prove the main result of this section, the following.

\textbf{12.3} \textit{For every integer $p \ge 0$ there exists $w \ge 0$ such that every graph with a vital $p$-linkage has path-width $\le w$.}

\textbf{Proof.} Choose $w'$ so that (1.1) holds with $w$ replaced by $w'$. Choose $n$ so that (2.6) holds with $k$ replaced by $w' + 1$. Let $w = \delm{w' + 3(w' + 1)2^{n-1}}$ $\fix{w' + 2^n(w' + 1)}$.\fixnote{fix:12.3-w}{The last sentence of the proof ends with ``$\le w' + 2^n(w' + 1) = w$''. The two values differ (for $n = 2$, they are $w' + 6(w' + 1)$ and $w' + 4(w' + 1)$). The second is the smaller, since $2^n < 3 \cdot 2^{n-1}$, so the proof also gives the result for the first.} We claim that $w$ satisfies (12.3). For let $G$ be a graph with a vital $p$-linkage $L$. By (1.1), $G$ has tree-width $\le w'$. By (11.2), there is a linked tree-decomposition $(T, W)$ of $G$, with width $\le w'$, such that every vertex of $T$ has degree at most 3. Let $N(T)$ be the set of vertices of $T$ with degree 3 in $T$, and let us choose $(T, W)$ with $|N(T)|$ minimum. \fix{By deleting edges from the graphs $W(t)$ as in the proof of (11.2), which changes neither $T$ nor the sets $W(e)$, we may also assume that no edge of $G$ belongs to $W(t)$ for two distinct $t \in V(T)$; so the pairs $(W(S), W(S'))$ below, for the two components $S, S'$ of $T$ minus an edge, are separations.}\fixnote{fix:12.3-trim}{Claims (1) and (2) below apply (12.1) and (2.6) to such pairs, so these must be separations (see the note in Section 11, page 43). The trimmed decomposition is still linked, of width $\le w'$, and has the same tree.}

\hspace*{12pt}(1) \textit{Let $f_1, f_2 \in E(T)$ with $|W(f_1)| = |W(f_2)| = k$ say, and let $T_0, T_1, T_2$ be the three components of}
